\section{The Chokepoint, the Record and the Literature}\label{sec:context}

\subsection{ASML and the record since 2019}\label{sec:background}

ASML is the only maker of extreme ultraviolet (EUV) lithography machines, on which the leading-edge logic of the most advanced US AI accelerators depends \citep{miller2022chip, asml2026annual}. China, which cannot buy them, buys immersion deep-ultraviolet (DUV) tools, and \citet{aei2026lithography} estimate that Chinese entities took 70 percent of ASML's immersion DUV sales in 2024. Exposing a layer several times, these tools print features that would otherwise need EUV. Customers in China accounted for 26.3 percent of ASML's net sales in 2023, 36.1 percent in 2024 and 29.1 percent in 2025 \citep{asml2026annual}. Servicing follows each sale. Installed-base management made up 25.1 percent of ASML's 2025 sales \citep{asml2026annual}, and several hundred ASML immersion tools in China need specialized maintenance about every six months, often from ASML staff on site \citep{aei2026lithography}.

Washington reaches the firm. At the end of 2025 ASML had 8,242 employees in the United States, 18.7 percent of its workforce \citep{asml2026annual}, and US subsidiaries including Cymer, the maker of its EUV light sources \citep{asml2026subsidiaries}. Under US export rules these entities and their US staff are ``U.S. persons'' \citep{ecfr2026usperson}, subject to US rules that restrict support for advanced-node fabs in China \citep{asml2026annual}. It also reaches the machines. An October 2023 rule set the de minimis threshold of US-origin content at zero for certain immersion tools bound for China \citep{ecfr7344, bis2023sme}. A rule effective on 2 December 2024 added a foreign direct product (FDP) rule and a de minimis rule specific to semiconductor manufacturing equipment (SME). Exporters in the 33 countries of its new Supplement No.~4 (25 EU members, the Netherlands among them, and Japan, but not South Korea or Taiwan) are excluded from the new licence requirements for items caught only by these two rules, because the listed governments ``have the authority to control key SME items of concern'' \citep{bis2024fdpr}. That status, which BIS can amend by rule, is what The Hague received (\cref{sec:coding} codes it).

Licences for ASML's shipments to China are a Dutch decision (\cref{sec:gatekeeper}). The Hague licensed an EUV sale to China in 2018 but, after US lobbying in 2019 that Reuters reported, did not renew the expired licence; ASML's new application was still pending in October 2020 \citep{nltimes2020pressured, kaag2020asml}. From 2023 each step followed a US move. National licensing of the most advanced immersion tools (NXT:2000i and later) took effect on 1 September 2023, after the US rules of October 2022 \citep{govnl2023controls}. Licensing of the older NXT:1970i and NXT:1980i moved from Washington to The Hague on 7 September 2024, after the US rule of October 2023 \citep{stcrt2024duv, asml2024license}. Controls on metrology and inspection systems took effect on 1 April 2025, after the December 2024 rule \citep{stcrt2025metrology}. On 15 November 2025 Commission Delegated Regulation (EU) 2025/2003 put the most advanced immersion tools on the EU list \citep{eu2025dualuse}, which ASML reported ``did not mean a change'': the licence decision stayed in The Hague \citep{asml2026annual}.

The US-facing direction is collective. EUV and the most advanced immersion tools ship to US fabs under the Union general export authorisation EU001, which a member state may withdraw from an exporter only for doubts about its ability to comply \citep{eu2021dualuse}. Only the NXT:1970i and NXT:1980i class, left on the Dutch national list, still need a Dutch licence to leave the EU \citep{stcrt2025_38653}. Turning the listed tools against Washington would take an EU decision (\cref{sec:council}).

Servicing in China remained permitted at least into mid-2026 \citep{zakaria2026match}. The MATCH Act, introduced in April 2026 and not enacted by late October, would ban immersion-tool sales anywhere in China and servicing at listed Chinese fabs, and would give allies 150 days to align before Washington acted alone \citep{baumgartner2026match, hr8170}. The Hague opposed the bill's extraterritorial reach \citep{tk2026match}; its trade minister said the bill ``could potentially put our national security in American hands'' \citep{nltimes2026irritated}.

Washington holds other instruments. The core of its guarantee to Europe is statutory: Section 1250A of the FY2024 National Defense Authorization Act bars withdrawal from the North Atlantic Treaty without two-thirds of the Senate or an Act of Congress \citep{usc1928f}. Posture is cheap to move. On 1 May 2026, after Chancellor Merz criticized US strategy on Iran, Defense Secretary Hegseth ordered about 5,000 troops out of Germany \citep{martinez2026germany}. Dependence runs both ways, though: on 30 March 2026 Spain, already refusing US use of jointly operated bases for the Iran war, closed its airspace to the US aircraft involved \citep{roque2026spain}. The AI Diffusion Rule of January 2025 put ten EU members in its favoured group and held the other seventeen to country allocations \citep{bis2025diffusion}, and BIS announced its rescission two days before compliance was due \citep{bis2025rescission}. In the August 2025 joint statement the EU said it intended to buy at least \$40 billion of US AI chips and to align its technology-security requirements with US ones, and Washington said it would then facilitate those exports \citep{whitehouse2025jointstatement}. The National Security Strategy of December 2025 holds out ``more favorable treatment on commercial matters, technology sharing, and defense procurement'' to countries that take more responsibility for their regional security and ``align their export controls with ours'' \citep{whitehouse2025nss}.

Beijing has built instruments against those who restrict it. Its extraterritorial rare-earth controls of October 2025 \citep{mayerbrown2025rareearth} were suspended until 10 November 2026 after the Busan summit of 30 October 2025 \citep{whitehouse2025busan, mofcom2025pause}. In September 2026 the two governments agreed to extend their truce to 10 January 2027, but we found no Chinese act extending the suspension itself \citep{mofcom2026extension}; \cref{sec:tests} asks what a lapse would bring. China's 2021 blocking rules let it penalize compliance with foreign extraterritorial measures \citep{squire2026blocking, geopolitechs2026blocking}, and two State Council decrees of April 2026 widened those powers to countermeasures against foreign firms that cut off normal trade with Chinese entities (Decree 834) and penalties for those who implement foreign extraterritorial measures (Decree 835) \citep{statecouncil2026supplychain, reedsmith2026decree}.

\subsection{What existing theories predict}\label{sec:literature}

\paragraph{Weaponized interdependence.} \citet{hirschman1945national} located influence in the losses a partner would suffer if trade stopped. \citet{farrell2019weaponized} moved the argument to networks and attached two conditions: a state must have ``physical or legal jurisdiction over hub nodes'' (p.~56) and the ``legal and regulatory institutions'' to exploit them (p.~57), and where hubs span jurisdictions the states involved must coordinate (pp.~56, 58). ASML meets the network condition, but Washington shares jurisdiction over it (\cref{sec:background}), and the European institutions that could turn it against the United States are collective (\cref{sec:council}). These works disagree about whether a hub yields lasting power. For Farrell and Newman the power comes from network position and the institutions behind it, and targets' responses come later \citep[76]{farrell2019weaponized}. \citet{gholz2021market} find that Japan absorbed China's 2010 rare-earth embargo through market adjustment and argue that cost measures that ignore market structure overstate a sender's power, and work on chokepoints as depleting assets sides with them: activation induces indigenization and rerouting \citep{tan2025driven, rajan2026eroding}. We let the exit premium drift and ask what a patron does while it lasts (\cref{sec:control}). \citet{wagner1988economic} supplies the discipline: asymmetric interdependence yields influence only when trading economic resources for political concessions beats bargaining over the economic gains alone, however asymmetric the relationship.

\paragraph{Alliance dependence and hierarchy.} \citet{snyder1984security, snyder1997alliance} distinguishes a protégé's need for its patron from the patron's own strategic interest in defending it; when that interest is well known, the stronger state ``cannot credibly threaten defection or realignment'' \citep[472--473]{snyder1984security}. \citet{morrow1991alliances} treats asymmetric alliances as trades of autonomy for security, and \citet{lake2009hierarchy} treats the patron's authority as relational, a social contract in which the patron provides order and protection and the subordinate complies. On those accounts compliance needs no fresh threat, a rival reading we take up in \cref{sec:predictions}.

The coercion literature disagrees about which instrument worked in the canonical case. \citet{gerzhoy2015alliance} credits conditional threats of military abandonment, paired with assurance, for stopping West Germany's nuclear ambitions. \citet{lanoszka2018atomic} finds alliance coercion hard and argues that an ally's economic and technological reliance on Washington is the stronger instrument. For \citet{pauly2024damned} the pairing is fragile, because the steps that make a threat credible undermine the assurance that compliance will not be punished; his assurance is owed to the target, while our cost of conditionality falls mainly on other allies (\cref{sec:conditionality}). \citet{wolford2017alliances} show that threats against a politically valuable ally can lack credibility, which is our channel of self-harming sanctions. A host that holds an asset its patron needs can collect from both sides. In 2009 Kyrgyzstan announced the closure of the US base at Manas alongside a Russian aid package, collected the first Russian payment, and then signed a new deal with Washington at a higher rent \citep[ch.~7]{cooley2012great}. Work on a patron restraining an ally's sensitive exports splits on motive: \citet{kroenig2009exporting} finds that sensitive nuclear transfers follow suppliers' strategic position and no support for commercial motives, while \citet{gray2012commercial} shows Bonn putting commercial rivalry first in its 1975 sale to Brazil.

\paragraph{Rival bidders and vote buying.} That a contested client earns what the losing bidder would pay is old. It is the one-voter case of sequential vote buying \citep{groseclose1996buying}; \citet[Prop.~2]{dekel2008vote} show that the winner pays the loser's value net of the client's preference, a price pinned down only in undominated strategies. \citet{bueno2016competition} show that a dominant aid donor captures the surplus of an aid-for-policy deal while a rival donor's entry raises what it must pay. Applied to ASML, this strand prices Dutch alignment at what China would pay for exit. The outside-option principle supplies the threshold: a threat to walk away matters only when the alternative beats continued bargaining \citep{shaked1984involuntary, sutton1986noncooperative, binmore1989outside}. In the two-hegemon extension of \citet[Supplement B.2]{clayton2026framework}, hegemons whose threats overlap bid transfers down to zero. \Cref{lem:bidding} restates these results; we claim no novelty for it. The effect of a rival is contested. Rival offers come when a protégé's relations with its patron sour \citep{yang2026major}, yet \citet{camboni2025spheres} find that China's growth can lead countries in the US sphere to offer Washington more. \citet{lacy2004theory} model sanctions as coercive linkage of an economic relationship to a political demand under two-sided incomplete information, so that imposed sanctions are a selected sample.

\paragraph{Slack.} \citet{clayton2026framework} show that a hegemon gains by pooling threats across relationships, bounded by the target's participation constraint, and \citet{clayton2026putting} state the principle we build on: ``Threats generate power to the extent that they generate slack in the participation constraint.'' Their pooled threats are enforcement threats that raise the target's inside option and are never spent; threats to cut off a target that refuses appear only in their Supplement B.3.1, where they are partly self-defeating. In \citet{clayton2026theory} a hegemon gains by committing to limit its own coercion.

Our patron's threats are cut-off threats aimed at an ally for whom a rival bids, each must be sequentially rational, and an executed threat is spent. As a result only conditional parts are capped by slack, conditional bites share one budget, and a patron that cannot commit to stop demanding obtains less (\cref{prop:slack}). Bidders in \citet{bueno2016competition} have carrots only. Ours have sticks and counter-threats, and China executes a counter-threat only to crowd out another demand (\cref{prop:counter}). Once the holder can propose and talks can rupture, the patron's dependence returns through agenda power, which the passive client of \citet{groseclose1996buying} and \citet{dekel2008vote} lacks (\cref{prop:master}). Where the hegemon of \citet{clayton2026theory} gains by binding itself, our patron chooses how conditional to make its guarantee, and it is the holder that may bind itself (\cref{prop:conditionality,rem:bind}).

\paragraph{Commitment and prevention.} Anticipated shifts in power create commitment problems that can make bargaining fail under complete information \citep{fearon1995rationalist, powell2006war}. Models of investment in future power divide over the response: uncertainty about the investment can itself produce preventive war \citep{debs2014known}, an imperfect monitor waits until the shift looks imminent \citep{bas2016dynamic}, and potential power can be bought off with concessions \citep{spaniel2019bargaining}; \cref{sec:preventive} sets our preventive demand for control against these models.
